\documentclass[final]{beamer}
\usepackage[size=a1,scale=1.15]{beamerposter}
\usepackage{tikz}
\usepackage{lmodern}
\usepackage{booktabs}
\usepackage{ragged2e}
\usetheme{default}
\definecolor{ink}{HTML}{17324D}
\definecolor{teal}{HTML}{0B7285}
\definecolor{mist}{HTML}{E8F1F5}
\setbeamercolor{normal text}{fg=ink,bg=white}
\setbeamercolor{block title}{fg=white,bg=ink}
\setbeamercolor{block body}{fg=ink,bg=mist}
\setbeamertemplate{navigation symbols}{}
\setbeamertemplate{block alerted begin}{\begin{beamercolorbox}[rounded=true,sep=1.2em]{block body alerted}\justifying}
\setbeamertemplate{block alerted end}{\end{beamercolorbox}}
\newcommand{\metric}[2]{\begin{beamercolorbox}[rounded=true,sep=0.7em,wd=\linewidth]{block body}\centering{\fontsize{34}{38}\selectfont\color{teal}\bfseries #1}\\[-2pt]{\large #2}\end{beamercolorbox}}
\title{A compact observation can reveal a system-level pattern}
\author{Maya Chen \and Theo Martins}
\institute{North Coast Systems Lab \quad\texttt{research@example.org}}
\begin{document}
\begin{frame}[t]
  \begin{beamercolorbox}[wd=\paperwidth,sep=1.4em,leftskip=2em,rightskip=2em]{block title}
    {\fontsize{42}{48}\selectfont\bfseries\inserttitle}\par\vspace{0.3em}
    {\Large\insertauthor}\hfill{\large\insertinstitute}
  \end{beamercolorbox}
  \vspace{1.1em}
  \begin{columns}[t,totalwidth=0.96\paperwidth]
    \begin{column}{0.31\paperwidth}
      \metric{37\%}{faster detection}
      \vspace{0.9em}
      \begin{block}{Question}
        Can a small, carefully chosen observation replace a long monitoring
        window when the underlying process changes in bursts?
      \end{block}
      \begin{block}{Approach}
        We combine an adaptive sampler with a lightweight change detector.
        The sampler spends its budget near transitions and the detector
        reports both the change point and its confidence.
        \begin{enumerate}\itemsep0.25em
          \item Calibrate on a quiet baseline.
          \item Sample densely when the signal bends.
          \item Confirm with an independent window.
        \end{enumerate}
      \end{block}
    \end{column}
    \begin{column}{0.36\paperwidth}
      \begin{alertblock}{Main finding}
        Adaptive sampling reaches the same decision quality with 43\% fewer
        observations, while preserving the false-alarm rate of the fixed
        schedule. The gain comes from spending attention at the boundary,
        not from changing the decision threshold.
      \end{alertblock}
      \vspace{0.8em}
      \begin{block}{Signal and decision boundary}
        \centering
        \begin{tikzpicture}[x=1cm,y=1cm]
          \draw[->,thick] (0,0)--(8.1,0) node[right]{time};
          \draw[->,thick] (0,0)--(0,4.2) node[above]{signal};
          \draw[teal,very thick] plot[smooth] coordinates{(0,1.0)(1,1.2)(2,1.1)(3,1.4)(4,2.0)(5,3.0)(6,3.1)(7,3.3)(8,3.2)};
          \draw[ink,dashed] (0,2.2)--(8,2.2) node[right]{threshold};
          \fill[teal!20] (3.7,0)--(5.5,0)--(5.5,4)--(3.7,4)--cycle;
          \node[teal,align=center] at (4.6,3.75){adaptive\\window};
        \end{tikzpicture}
      \end{block}
      \begin{block}{Evaluation}
        \centering
        \begin{tabular}{lrrr}
          \toprule
          Schedule & Recall & False alarms & Samples\\\midrule
          Fixed & 0.81 & 0.06 & 1200\\
          Adaptive & 0.82 & 0.06 & 684\\
          \bottomrule
        \end{tabular}
      \end{block}
    \end{column}
    \begin{column}{0.25\paperwidth}
      \metric{0.06}{false-alarm rate}
      \vspace{0.9em}
      \begin{block}{What to remember}
        \begin{itemize}\itemsep0.45em
          \item Concentrate measurements around transitions.
          \item Keep a confirmation window to protect precision.
          \item Report the budget and uncertainty together.
        \end{itemize}
      \end{block}
      \begin{block}{Next step}
        Test the sampler on a second site with a different noise profile and
        pre-register the alert budget before collecting data.
      \end{block}
      \vfill
      {\small\textbf{Code/data:} example.org/project\\
      \textbf{Acknowledgements:} Coastal Data Initiative}
    \end{column}
  \end{columns}
\end{frame}
\end{document}
